\section{\texorpdfstring{\texttt{spec4pbt}}{spec4pbt}: Trustworthy PBT toolchain for SpecTec Specifications}
\label{sec:instantiation}

\subsection{Overview}
\label{sec:inputs}
\label{sec:overview}

\texttt{spec4pbt} is a tool for SpecTec specifications that addresses the three challenges introduced in Section~\ref{sec:three-propositions} (Figure~\ref{fig:overview}).

\begin{figure}[H]
  \centering
  \resizebox{\linewidth}{!}{% One Lean context contains both mechanically derived and agent-authored code.
% Orange arrows: agent construction. Blue arrows: Lean proof checking only.
% Dashed arrows: mechanical translation or derivation. Solid black arrows: testing data flow.
\begin{tikzpicture}[
  font=\fontsize{9}{11}\selectfont,
  box/.style={draw,rounded corners=2pt,align=center,minimum height=8mm,inner sep=3pt},
  input/.style={box,rounded corners=0pt,thick,fill=white,text width=2.85cm,minimum height=1cm},
  doc/.style={chamfered rectangle,chamfered rectangle corners=north east,
    chamfered rectangle xsep=1.8mm,draw,thick,align=center,minimum height=9mm,inner sep=3pt},
  leandoc/.style={doc,draw=blue!60!black,fill=blue!9},
  parsed/.style={box,draw=black!80,thick,fill=white,minimum width=1.8cm},
  edge/.style={-{Stealth[length=1.8mm]},thick},
  agentarrow/.style={edge,very thick,draw=orange!75!black},
  leanarrow/.style={edge,very thick,draw=blue!60!black},
  scriptarrow/.style={edge,dashed,draw=black!75},
  lbl/.style={font=\fontsize{8}{9.5}\selectfont,align=center,inner sep=2pt,text=black},
  agentlbl/.style={lbl,text=orange!75!black,font=\fontsize{8}{9.5}\selectfont\bfseries},
  leanlbl/.style={lbl,text=blue!60!black,font=\fontsize{8}{9.5}\selectfont\bfseries},
  heading/.style={font=\fontsize{9}{11}\selectfont\bfseries,anchor=south west,text=black}
]
  % The frame denotes the language of the artifacts, not who creates them.
  \begin{scope}[on background layer]
    \draw[draw=black!55,thick,fill=black!2] (-2.3,-2.55) rectangle (1.1,3.75);
    \draw[rounded corners=4pt,draw=blue!40,fill=blue!4] (3.9,-2.55) rectangle (12.75,3.75);
  \end{scope}
  \node[heading] at (-2.1,3.25) {Inputs};
  \node[heading,text=blue!60!black] at (4.05,3.25) {Lean definitions, propositions, and proofs};

  % Source inputs remain plain rectangular boxes.
  \node[input] (sourcespec) at (-0.6,2.5) {SpecTec\\specification};
  \node[input] (sourceproperty) at (-0.6,0.35) {Property in\\SpecTec-like notation};
  \node[input] (implementation) at (-0.6,-1.9) {Implementation};
  \node[parsed] (parsedproperty) at (2.85,0.35) {Parsed\\property};
  \draw[scriptarrow] (sourceproperty.east) -- node[lbl,above] {parse} (parsedproperty.west);

  \node[leandoc,minimum width=3cm] (specification) at (6.0,2.5) {Specification in Lean};
  \node[leandoc,minimum height=1.15cm] (property) at (6.0,0.35)
    {Property in Lean\\[2pt]$\mathit{Precondition}(x)\Rightarrow\mathit{Goal}(x)$};
  \node[leandoc,minimum width=3cm] (generatorprop) at (10.5,0.35)
    {Generator sound and\\complete for\\$\mathit{Precondition}(x)$};
  \draw[scriptarrow] (sourcespec.east) -- node[lbl,above] {mechanically translate} (specification.west);

  % All three outputs come from the parsed source property, not the Lean proposition.
  \draw[scriptarrow] (parsedproperty.east) -- (property.west);
  \draw[scriptarrow] (parsedproperty.north) -- (2.85,1.55)
    -- node[lbl,above,pos=0.72] {from Precondition} (10.5,1.55) -- (generatorprop.north);

  \node[leandoc,minimum width=1.65cm] (propertyproof) at (6.0,-1.9) {Property\\proof};
  \node[leandoc,minimum width=1.65cm] (generatorproof) at (9.4,-1.9) {Generator\\proofs};
  \node[leandoc,minimum width=1.65cm] (generator) at (11.7,-1.9) {Input\\generator};
  \draw[agentarrow] ([xshift=-2.5mm]property.south) -- node[agentlbl,left] {agent\\proves} ([xshift=-2.5mm]propertyproof.north);
  \draw[leanarrow] ([xshift=2.5mm]propertyproof.north) -- node[leanlbl,right] {Lean\\checks} ([xshift=2.5mm]property.south);
  % The agent constructs both code and proofs; only proofs establish the propositions.
  \draw[very thick,draw=orange!75!black] (generatorprop.south)
    -- node[agentlbl,right] {agent\\constructs} (10.5,-1.0);
  \draw[agentarrow] (10.5,-1.0) -| ([xshift=2.5mm]generatorproof.north);
  \draw[agentarrow] (10.5,-1.0) -| (generator.north);
  \draw[leanarrow] ([xshift=-2.5mm]generatorproof.north)
    -- node[leanlbl,left] {Lean\\checks} ([xshift=-2.5mm]generatorproof.north |- generatorprop.south);

  % The checker receives the generated test inputs and the supplied implementation.
  \node[doc,draw=black,fill=white,minimum width=2.6cm] (checker) at (11.7,-4.0)
    {Checker\\for $\mathit{Goal}(x)$};
  \draw[scriptarrow] (parsedproperty.south) -- (2.85,-4.0)
    -- node[lbl,above,pos=0.7] {from Goal} (checker.west);
  \node[box,minimum width=1.7cm] (verdict) at (14.4,-4.0) {Pass / fail};
  \draw[edge] (generator.south) -- node[lbl,right] {sample\\test inputs} (checker.north);
  \draw[edge] (implementation.east) -- (1.45,-1.9) -- (1.45,-5.0) -- (11.7,-5.0) -- (checker.south);
  \draw[edge] (checker.east) -- (verdict.west);

  % Arrow styles describe operations independently of the Lean artifact frame.
  \draw[scriptarrow] (-2.1,-5.8) -- (-1.3,-5.8);
  \node[lbl,anchor=west] at (-1.2,-5.8) {Mechanical translation / derivation (\texttt{spec4pbt})};
  \draw[edge] (7.9,-5.8) -- (8.7,-5.8);
  \node[lbl,anchor=west] at (8.8,-5.8) {Runtime data flow};
\end{tikzpicture}
}
  \caption{Overview of \texttt{spec4pbt}. Orange arrows are agent construction and blue arrows are Lean proof checking.}
  \label{fig:overview}
\end{figure}

\subsection{SpecTec and Its Translation to Lean}
\label{sec:spectec}
\label{sec:translation}

% Avoid stretched spacing around examples in the outline.
\raggedbottom

% Topic sentence for the specification framework used by the instantiation.
SpecTec is a domain-specific language for mechanizing language specifications, introduced for the WebAssembly standard~\cite{youn24spectec} and later adapted to P4~\cite{lee26p4spectec}.
It expresses syntax and semantics as structured definitions.


An excerpt of the P4 implicit-cast rules shows the main SpecTec constructs (Figure~\ref{fig:ordered-translation}(a)).
The translation maps syntax definitions to inductive types and rules to clauses that \texttt{firstSuccess} tries in order (Figure~\ref{fig:ordered-translation}(b)).

\begin{figure}[H]
  \centering
  % (a) Source: p4-spectec/spec/5-typing/5.04-typing-relation.watsup (castResult,
% Cast_impl_neq) and 5.05.2-typing-casting.watsup (rules). typeIR is flattened,
% prose hints and rulegroups are dropped, and the remaining Cast_impl_neq rules
% are elided.
\begin{examplerow}
\begin{minipage}[t]{0.5\linewidth}
\vspace{0pt}
\begin{lstlisting}[style=spectecexample,style=examplecolumn]
syntax typeIR =
  | INT
  | INT `< nat `>
  | BIT `< nat `>
  | ...

syntax castResult =
  | OK typeIR*
  | FAIL
\end{lstlisting}
\end{minipage}%
\begin{minipage}[t]{0.5\linewidth}
\vspace{0pt}
\begin{lstlisting}[style=spectecexample,style=examplecolumn]
relation Cast_impl_neq:
  typeIR -> typeIR : castResult
  hint(input %0 %1)

rule Cast_impl_neq/fixInt:
  INT -> INT `< w `> : OK ([ INT `< w `> ])
rule Cast_impl_neq/fixBit:
  INT -> BIT `< w `> : OK ([ BIT `< w `> ])
;; ... other rules
rule Cast_impl_neq/fail:
  typeIR_a -> typeIR_b : FAIL
  -- otherwise
\end{lstlisting}
\end{minipage}
\end{examplerow}
\examplepanelcaption{(a) P4-SpecTec excerpt.}

% (b) Generated by the ordered backend from (a). Names are shortened, pattern
% guards are condensed, the fail clause is inlined, and noncomputable
% annotations are omitted.
\begin{examplerow}
\begin{minipage}[t]{0.5\linewidth}
\vspace{0pt}
\begin{lstlisting}[style=leanexample,style=examplecolumn]
inductive typeIR where
  | INT : typeIR
  | «INT `< % `>» : Nat → typeIR
  | «BIT `< % `>» : Nat → typeIR
  -- ...

inductive castResult where
  | OK : List typeIR → castResult
  | FAIL : castResult
\end{lstlisting}
\end{minipage}%
\begin{minipage}[t]{0.5\linewidth}
\vspace{0pt}
\begin{lstlisting}[style=leanexample,style=examplecolumn]
def Cast_impl_neq_0 (t1 t2 : typeIR) :
    Plan Request :=
  .bind (.guard (t1 = .INT)) (fun _ =>
  match t2 with
  | .«INT `< % `>» w =>
    .pure (.OK [.«INT `< % `>» w])
  | _ => .unmatch)

def Cast_impl_neq_plan (t1 t2 : typeIR) :
    Plan Request :=
  firstSuccess
    [fun _ => Cast_impl_neq_0 t1 t2,
     fun _ => Cast_impl_neq_1 t1 t2,
     fun _ => .pure .FAIL] () -- otherwise

def Cast_impl_neq (t1 t2 : typeIR)
    (r : castResult) : Prop :=
  ∃ n, Eval program
    (.call (.Cast_impl_neq t1 t2)) 0
    (.success r) n
\end{lstlisting}
\end{minipage}
\end{examplerow}
\examplepanelcaption{(b) Ordered Lean translation.}

  \caption{\texttt{Cast\_impl\_neq} takes two types as inputs and returns \texttt{FAIL} when no other rule applies (abbreviated). Generated names are shortened and pattern guards are condensed.}
  \Description{The top panel shows syntax definitions for typeIR and castResult, the Cast_impl_neq relation with two input positions, two rules casting INT to fixed-width integer types, and a fail rule guarded by otherwise. The bottom panel shows typeIR and castResult as inductive types with constructors named after the SpecTec notation, the fixInt rule as a clause over the two inputs that returns OK or unmatches, firstSuccess trying the clauses in order with FAIL last, and the relation defined as successful evaluation of the call.}
  \label{fig:ordered-translation}
\end{figure}

The translation's semantic correspondence with SpecTec requires separate validation, which we perform in Section~\ref{sec:validation}.

\subsection{Formalizing and Proving Properties}
\label{sec:propose}
\label{sec:derived-propositions}

We express a property in a restricted SpecTec-like syntax with explicit inputs, ordered premises, and a conclusion, from which scripts algorithmically derive Lean propositions and an executable checker.
We focus on the P4 property that a type-checked program instantiates successfully (Figure~\ref{fig:instantiation-property}).
\textcolor{red}{(Note: I ultimately intend a negated property such as ill-typed rejection; this example is selected to show how relation with outputs is translated.)}

\begin{figure}[H]
  \centering
  \begin{minipage}[b]{0.485\linewidth}
\begin{lstlisting}[style=spectecexample]
relation Program_ok:
  |- p4program : p4programIR
  hint(input %0)

relation Program_inst:
  |- p4program : globalInstLayer store
  hint(input %0)

property instantiation_progress:
  (p : p4program)
  Program_inst: |- p : l s
  -- Program_ok: |- p : ir
\end{lstlisting}
\examplepanelcaption{(a) Specification relations and SpecTec-like property.}
\end{minipage}\hfill
\begin{minipage}[b]{0.485\linewidth}
\begin{lstlisting}[style=leanexample]
structure Input where
  p : p4program

def Premise (input : Input) : Prop :=
  ∃ ir, Program_ok input.p ir

def Goal (input : Input) : Prop :=
  ∃ l s, Program_inst input.p l s

def property : Prop :=
  ∀ input : Input,
    Premise input → Goal input
\end{lstlisting}
\examplepanelcaption{(b) Derived Lean property contract.}
\end{minipage}

  \caption{The P4 instantiation property and its derived Lean contract (abbreviated); the property proof remains open.}
  \Description{The left panel shows the Program_ok and Program_inst relation signatures and the instantiation_progress property. The right panel shows the derived Input, Premise, Goal, and property definitions in Lean.}
  \label{fig:instantiation-property}
\end{figure}

\subsection{Constructing and Verifying Generators}
\label{sec:generators}

\subsubsection{The Overall Construction Problem}
\label{sec:construction}

The task is to construct a generator for the property's precondition and prove its soundness and completeness.
We adopt Palamedes's method of constructing a generator by proving the existence of one satisfying soundness and completeness~\cite{goldstein26palamedes}.

\begin{figure}[H]
\begin{lstlisting}[style=leanexample]
abbrev generatorValue := Gen Input

def generatorSound (g : generatorValue) : Prop :=
  ∀ input, Gen.support g input → Premise input

def generatorComplete (g : generatorValue) : Prop :=
  ∀ input, Premise input → Gen.support g input
\end{lstlisting}
\caption{The derived generator contract: a generator is sound if every input it can produce satisfies \texttt{Premise}, and complete if it can produce every input satisfying \texttt{Premise}.}
\Description{Lean definitions of generatorValue as a generator of Input, generatorSound stating that every input in the generator's support satisfies Premise, and generatorComplete stating that every input satisfying Premise is in the generator's support.}
\label{fig:generator-contract}
\end{figure}

\subsubsection{Decomposition into Smaller Tasks}
\label{sec:decomposition}

We can delegate generator construction and proof entirely to an AI agent using Lean feedback, but solving the whole problem at once is difficult.
We therefore explore decomposing generator construction and proof into smaller tasks, then composing and verifying their results against the original soundness and completeness propositions.

\subsection{Testing External Implementations}
\label{sec:test}

We exploit the executability of P4-SpecTec to derive a checker from the property's goal.

To test an external implementation, the user identifies the specification relations whose behavior is externally observable through that implementation.
Along with the implementation, the user supplies an adapter for each such relation, specifying how to invoke the implementation on the relation's inputs and map its observed outcomes to the relation's results.

From the property's goal, \texttt{spec4pbt} collects the required external relations and binds them to the supplied adapters.
It then algorithmically runs the implementation on generated inputs and evaluates the goal using the observed outcomes to produce pass/fail verdicts.

\flushbottom
